\section{Расчёт освещённости поверхностей}
\label{sec:lighting}

Модель освещения "--- функция, по которой вычисляется интенсивность света, отражённого к наблюдателю~\cite[с.~380--385]{rogers1989}. Далее $\vec{N}$ "--- единичная нормаль к поверхности, $\vec{V}$ "--- единичный вектор направления на наблюдателя, $C_{\text{п}}$ "--- цвет поверхности, $\circ$ "--- покомпонентное произведение цветов. Модель Ламберта учитывает только диффузное отражение и бликов не даёт~\cite[с.~380]{rogers1989}, поэтому требованию задания не удовлетворяет. Фоновая составляющая~\cite[с.~381]{rogers1989} не вводится: в ночной сцене её роль играет свечение сияния.
% (подраздел~\ref{sec:aurora-lighting}).

\subsection{Модель Фонга}

Модель Фонга~\cite[с.~382--385]{rogers1989} дополняет диффузное отражение по закону Ламберта зеркальной составляющей "--- бликом, тем более ярким, чем ближе $\vec{V}$ к направлению зеркального отражения
\begin{equation}
	\vec{R} = 2 (\vec{N} \cdot \vec{L})\, \vec{N} - \vec{L},
	\label{eq:reflect}
\end{equation}
и описывается формулой
\begin{equation}
	I = I_L\, C_L \circ \left( C_{\text{п}}\, k_d \max(0,\ \vec{N} \cdot \vec{L}) + k_s \max(0,\ \vec{R} \cdot \vec{V})^n \right).
	\label{eq:phong}
\end{equation}
Малые значения $n$ дают широкий блик матовых поверхностей~\cite[с.~383]{rogers1989}; цвет блика совпадает с цветом источника~\cite[с.~384]{rogers1989}.

\subsection{Модель Блинна "--- Фонга}

В модели Блинна "--- Фонга вместо $\vec{R} \cdot \vec{V}$ используется $\vec{N} \cdot \vec{H}$, где $\vec{H}$ "--- вектор полупути, направленный по биссектрисе угла между $\vec{L}$ и $\vec{V}$~\cite[с.~403--404]{rogers1989}, \cite{hearn2005}:
\begin{equation}
	\vec{H} = \frac{\vec{L} + \vec{V}}{|\vec{L} + \vec{V}|},
	\label{eq:halfway}
\end{equation}
\begin{equation}
	I = I_L\, C_L \circ \left( C_{\text{п}}\, k_d \max(0,\ \vec{N} \cdot \vec{L}) + k_s \max(0,\ \vec{N} \cdot \vec{H})^n \right).
	\label{eq:blinn-phong}
\end{equation}
Для освещённой и видимой точки $\vec{N} \cdot \vec{L} > 0$ и $\vec{N} \cdot \vec{V} > 0$, поэтому $\vec{N} \cdot \vec{H} > 0$, тогда как $\vec{R} \cdot \vec{V}$ может быть отрицательным, и широкий блик модели Фонга обрывается резким краем.

\subsection{Выбор модели освещения}

Число операций в таблице~\ref{tab:lighting} подсчитано по формулам~\eqref{eq:reflect}--\eqref{eq:blinn-phong} для одной точки и одного источника без общего множителя $I_L C_L$.

\begin{table}[H]
	\caption{Сравнение моделей освещения}
	\label{tab:lighting}
	\begin{tabular}{|L{0.44}|C{0.27}|C{0.27}|}
		\hline
		Критерий & Фонг & Блинн "--- Фонг \\
		\hline
		Умножений / сложений & 12 / 8 & 11 / 10 \\
		\hline
		Прочие операции & степень & корень, 3 деления, степень \\
		\hline
		Обрыв блика & возможен & невозможен \\
		\hline
	\end{tabular}
\end{table}

Модель Блинна "--- Фонга дороже на одно извлечение корня и три деления, но не даёт обрыва блика на матовых поверхностях, поэтому выбрана она.